\chapter{Expander Graphs}

In Chapter~33, we searched for bottlenecks to partition graphs into clusters. An expander is the opposite extreme: a graph with no bottlenecks anywhere, where every cut is large and information spreads quickly.

A good road network has a peculiar property. No matter how you try to split the cities into two halves, you always have to cut a lot of roads. There is no narrow waist, no bottleneck you can sever with a few cuts to isolate a large region. Every set of cities, unless it is almost everybody or almost nobody, has a thick boundary leaking out to the rest. Networks like this are robust: knock out a handful of links and the graph stays connected; start a rumor anywhere and it floods the whole network in a logarithmic number of rounds; run a random walk and it forgets where it started almost immediately.

These are \emph{expander graphs}, and they are among the most useful objects in theoretical computer science. They are sparse---only $\Oh(n)$ edges---yet they behave, in every way that matters for routing, mixing, and information flow, like the complete graph on $n$ vertices. This combination of dense-graph connectivity with a sparse graph's edge count is why they appear in so many places: in fault-tolerant networks, in derandomization, in error-correcting codes, in the analysis of Markov chains, and behind nearly optimal algorithms for cuts and flows.

This chapter develops expanders from three complementary angles that prove quantitatively equivalent. The \emph{combinatorial} angle measures expansion by counting boundary edges. The \emph{spectral} angle measures it by the gap between the first two Laplacian eigenvalues. The \emph{probabilistic} angle shows that a random sparse regular graph is, with high probability, a strong expander. We tie the combinatorial and spectral views together with Cheeger's inequality and the Expander Mixing Lemma, sketch how to build expanders explicitly (the genuinely difficult step), and then turn to applications. Throughout we build on the Laplacian and spectral tools of Chapters~31--33.

\chapterline

\section{Two ways to measure a bottleneck}
\label{ch34:two-measures}

Fix a graph $\Gph=(V,E)$ with $n=\abs{V}$ vertices. For most of this chapter we take $\Gph$ to be $d$-regular: every vertex has exactly $d$ neighbors. Regularity is not strictly necessary, but it keeps the algebra clean, the eigenvalues comparable, and the constructions natural. The fundamental question is: \emph{how hard is it to cut $\Gph$ into two substantial pieces?}

\begin{intuitionbox}
Picture a cycle $C_n$ versus the complete graph $K_n$. To split the cycle in half you snip exactly two edges; the cycle has a thin waist everywhere. To split $K_n$ in half you must cut roughly $n^2/4$ edges---a quarter of all pairs. The cycle is the worst expander; the complete graph is the best. An expander is a sparse graph (like the cycle, $\Oh(n)$ edges) that is hard to cut (like $K_n$). This chapter studies the tension between ``few edges'' and ``hard to cut.''
\end{intuitionbox}

For a set $S\subseteq V$ write $\bar S = V\setminus S$. The \emph{edge boundary} is
\[
  \partial S \;=\; \set{(u,v)\in E : u\in S,\ v\in\bar S},
\]
the set of edges with exactly one endpoint in $S$. Cutting $S$ away from the rest costs $\abs{\partial S}$ edges.

\begin{definition}[Edge expansion / Cheeger constant]
\label{ch34:def-edge-expansion}
The \emph{edge expansion} (or \emph{conductance}, or \emph{Cheeger constant}) of $\Gph$ is
\[
  h(\Gph)\;=\;\min_{\substack{S\subseteq V\\ 0<\abs{S}\le n/2}}\ \frac{\abs{\partial S}}{\abs{S}}.
\]
\end{definition}

The restriction $\abs{S}\le n/2$ matters: we normalize by the size of the \emph{smaller} side. Without it, taking $S = \emptyset$ or $S = V$ would produce an uninformative ratio of $0$. Reading the definition aloud: $h(\Gph)$ is large exactly when \emph{every} set of size at most $n/2$ has many edges escaping it. A large $h$ means no bottleneck.

\begin{definition}[Combinatorial expander]
\label{ch34:def-expander}
A family $\set{\Gph_n}$ of $d$-regular graphs (with $n\to\infty$, $d$ fixed) is a \emph{family of expanders} if there is a constant $\eps>0$ with $h(\Gph_n)\ge \eps$ for all $n$.
\end{definition}

The word ``family'' is essential. A single finite graph always has \emph{some} expansion; the content is that expansion stays bounded \emph{away from zero} as the graph grows, while the degree $d$---and hence the edge count---stays constant. The surprise is that arbitrarily large graphs with only $dn/2$ edges can still guarantee that a fixed positive fraction of any set's volume leaks out.

\begin{focusbox}
Three normalizations float around the literature, and conflating them is the classic source of off-by-a-factor confusion.
\begin{itemize}
  \item \textbf{Edge expansion} $h(\Gph)=\min \abs{\partial S}/\abs{S}$ (boundary per vertex; what Definition~\ref{ch34:def-edge-expansion} uses).
  \item \textbf{Conductance} $\phi(\Gph)=\min \abs{\partial S}/(d\abs{S})$ for $\abs{S}\le n/2$ (boundary per unit \emph{volume}; volume $=$ sum of degrees). For $d$-regular graphs $\phi = h/d$.
  \item \textbf{Vertex expansion} counts boundary \emph{vertices} rather than edges (below).
\end{itemize}
We use $h$ for combinatorics and $\phi$ when we want a quantity in $[0,1]$ that matches random-walk mixing. Always check which one a source means.
\end{focusbox}

There is a second, often stronger, notion that counts vertices instead of edges. Let $N(S)$ be the set of vertices outside $S$ adjacent to something in $S$, the \emph{vertex boundary}.

\begin{definition}[Vertex expansion]
\label{ch34:def-vertex-expansion}
$\Gph$ is a \emph{$(K,A)$ vertex expander} if every set $S$ with $\abs S\le K$ satisfies $\abs{N(S)}\ge A\abs{S}$.
\end{definition}

Vertex expansion with $A>1$ says small sets have \emph{more} outside neighbors than they have members---a strict ``growth'' condition. It is what bipartite expanders for error-correcting codes will need (Section~\ref{ch34:codes}). Edge and vertex expansion are related but not identical; a vertex of high degree can have a big edge boundary while touching few distinct vertices. For bounded degree $d$ they sandwich each other within constant factors: $\abs{\partial S}/d \le \abs{N(S)} \le \abs{\partial S}$, so $h/d \le$ (vertex expansion of small sets) $\le h$.

\begin{figure}[H]
\centering
\begin{tikzpicture}[
  u/.style={circle,draw=sourceblue,fill=sourceblue!15,thick,minimum size=6mm,font=\scriptsize\bfseries},
  v/.style={circle,draw=targetgreen,fill=targetgreen!15,thick,minimum size=6mm,font=\scriptsize\bfseries},
  inactive/.style={circle,draw=gray!50,fill=gray!10,thick,minimum size=6mm,font=\scriptsize,text=gray},
  edgeS/.style={draw=sourceblue!80,thick},
  edgeOther/.style={draw=gray!30,thin}]

  % Background panels
  \draw[rounded corners=6pt,fill=sourceblue!5,draw=sourceblue!25,thick] (-0.6,-2.9) rectangle (1.0,2.95);
  \node[font=\footnotesize\bfseries,text=sourceblue] at (0.2,3.25) {Left Side $L$};

  \draw[rounded corners=6pt,fill=targetgreen!5,draw=targetgreen!25,thick] (3.8,-2.9) rectangle (5.4,2.95);
  \node[font=\footnotesize\bfseries,text=targetgreen] at (4.6,3.25) {Right Side $R$};

  % Left nodes
  \node[u] (u1) at (0.2,2.2) {$u_1$};
  \node[u] (u2) at (0.2,1.1) {$u_2$};
  \node[u] (u3) at (0.2,0.0) {$u_3$};
  \node[inactive] (u4) at (0.2,-1.1) {$u_4$};
  \node[inactive] (u5) at (0.2,-2.2) {$u_5$};

  % Right nodes
  \node[v] (v1) at (4.6,2.3) {$v_1$};
  \node[v] (v2) at (4.6,1.35) {$v_2$};
  \node[v] (v3) at (4.6,0.4) {$v_3$};
  \node[v] (v4) at (4.6,-0.55) {$v_4$};
  \node[v] (v5) at (4.6,-1.5) {$v_5$};
  \node[inactive] (v6) at (4.6,-2.45) {$v_6$};

  % Bounding box for S (label placed outside outer box to avoid border collision)
  \draw[sourceblue,dashed,thick,rounded corners=4pt] (-0.3,-0.45) rectangle (0.7,2.65);
  \node[font=\scriptsize\bfseries,text=sourceblue,left=4pt] at (-0.6,1.1) {Subset $S$ ($\abs{S}=3$)};

  % Bounding box for N(S) (label placed outside outer box to avoid border collision)
  \draw[targetgreen,dashed,thick,rounded corners=4pt] (4.1,-1.9) rectangle (5.1,2.75);
  \node[font=\scriptsize\bfseries,text=targetgreen,right=4pt,align=left] at (5.4,0.4) {Neighborhood $N(S)$\\($\abs{N(S)}=5 \ge A\abs{S}$)};

  % Expanding fan-out edges from S to N(S)
  \draw[edgeS] (u1) -- (v1); \draw[edgeS] (u1) -- (v2); \draw[edgeS] (u1) -- (v3);
  \draw[edgeS] (u2) -- (v2); \draw[edgeS] (u2) -- (v4); \draw[edgeS] (u2) -- (v5);
  \draw[edgeS] (u3) -- (v1); \draw[edgeS] (u3) -- (v3); \draw[edgeS] (u3) -- (v5);

  % Inactive edges from L \ S
  \draw[edgeOther] (u4) -- (v4); \draw[edgeOther] (u4) -- (v6);
  \draw[edgeOther] (u5) -- (v5); \draw[edgeOther] (u5) -- (v6);

  % Fan-out annotation placed below all edges
  \node[font=\scriptsize\itshape,text=edgegray,align=center] at (2.4,-3.35)
    {High degree fan-out ($d=3$): few overlapping neighbors $\implies$ rapid boundary expansion};
\end{tikzpicture}
\caption{Bipartite expander neighbor fan-out. A subset $S$ of $3$ variable nodes on the left connects to
$\abs{N(S)} = 5$ distinct constraint nodes on the right via non-colliding edges. In a lossless expander,
every small set $S$ achieves expansion $\abs{N(S)} \ge (1-\eps)d\abs{S}$, guaranteeing that most neighbors
touch $S$ at a unique edge.}
\label{ch34:fig-fanout}
\end{figure}

\begin{labbox}
Take the $3$-cube $Q_3$: vertices are bit strings $000,\dots,111$, edges connect strings differing in one bit. It is $3$-regular on $8$ vertices. Pick $S=\set{000,001,010,011}$ (all strings starting with $0$). Its boundary edges are exactly the four ``flip the first bit'' edges $000$--$100$, $001$--$101$, $010$--$110$, $011$--$111$. So $\abs{\partial S}=4$, $\abs S=4$, ratio $1$. Now try $S=\set{000}$: boundary $=3$ (its three neighbors), ratio $3$. Try $S=\set{000,111}$ (an antipodal pair, not adjacent): each contributes $3$ boundary edges, none shared, ratio $6/2=3$. The half-cube split is the cheapest, so $h(Q_3)=1$. Verify by checking a couple more $4$-sets that none beats ratio $1$.
\end{labbox}

\begin{figure}[H]
\centering
\begin{tikzpicture}[scale=1.1,
  v/.style={circle,draw=nodeblue,fill=nodeblue!12,inner sep=1.6pt,font=\scriptsize},
  cut/.style={trapred,line width=1.4pt}]
  \begin{scope}[shift={(-3.4,0)}]
    \foreach \i in {0,...,7}{
      \node[v] (c\i) at ({90-45*\i}:1.2) {};
    }
    \foreach \i [evaluate=\i as \j using {int(mod(\i+1,8))}] in {0,...,7}{
      \draw (c\i)--(c\j);
    }
    \draw[cut] (c1)--(c2);
    \draw[cut] (c5)--(c6);
    \node at (0,-1.7) {\scriptsize cycle: cut $2$ edges};
  \end{scope}
  \begin{scope}[shift={(3.4,0)}]
    \foreach \i in {0,...,7}{
      \node[v] (k\i) at ({90-45*\i}:1.2) {};
    }
    \foreach \i in {0,...,7}{
      \foreach \j in {0,...,7}{
        \ifnum\i<\j \draw[edgegray!50,line width=0.3pt] (k\i)--(k\j);\fi
      }
    }
    \node at (0,-1.7) {\scriptsize $K_8$: cut $\sim 16$ edges};
  \end{scope}
  \node at (0,0) {\large $\Longleftrightarrow$};
  \node at (0,0.45) {\scriptsize expander};
  \node at (0,-0.45) {\scriptsize lives here};
\end{tikzpicture}
\caption{The two extremes. A cycle has a thin waist (small $h$); a complete graph has a thick waist (large $h$). An expander wants the cut-resistance of $K_n$ at the edge-budget of the cycle.}
\label{ch34:fig-extremes}
\end{figure}

\section{The spectral viewpoint and Cheeger's inequality}
\label{ch34:spectral}

Computing $h(\Gph)$ exactly means minimizing over all $2^{n-1}$ subsets; the associated decision problem is NP-hard (it is essentially sparsest cut). The way around this, and the central idea of the chapter, is that a single eigenvalue of the Laplacian \emph{estimates} $h$ to within a quadratic factor, and that eigenvalue is computable in polynomial time.

Recall from Chapter~31 the Laplacian $\Lap=\Adjm_D-\Adjm$, where $\Adjm_D$ is the diagonal degree matrix and $\Adjm$ the adjacency matrix. For a $d$-regular graph $\Adjm_D=dI$, so $\Lap=dI-\Adjm$. The key identity is
\[
  x^\top \Lap x \;=\; \sum_{(u,v)\in E}\bigl(x_u-x_v\bigr)^2,
\]
the sum of squared differences across edges---a measure of how non-smooth $x$ is along the graph. $\Lap$ is symmetric positive semidefinite, so its eigenvalues are real and nonnegative:
\[
  0=\lambda_1\le \lambda_2\le\cdots\le\lambda_n.
\]
The all-ones vector $\mathbf 1$ has $\Lap\mathbf 1=0$, giving $\lambda_1=0$. The multiplicity of $0$ equals the number of connected components, so $\lambda_2>0$ iff $\Gph$ is connected. The quantity $\lambda_2$ is the \emph{spectral gap} (for regular graphs; the algebraic connectivity or Fiedler value of Chapters~31 and~32). It is the single number that controls expansion.

\begin{intuitionbox}
Why should an eigenvalue know about cuts? Suppose $\Gph$ had a sparse cut $(S,\bar S)$. The indicator-like vector that is $+1/\abs S$ on $S$ and $-1/\abs{\bar S}$ on $\bar S$, made orthogonal to $\mathbf 1$, is nearly constant inside each side and only ``pays'' $x^\top\Lap x$ on the few boundary edges. So it is a cheap test vector, and by the variational (Courant--Fischer) characterization
\[
  \lambda_2=\min_{x\perp\mathbf 1,\ x\ne 0}\frac{x^\top\Lap x}{x^\top x},
\]
a cheap test vector forces $\lambda_2$ to be small. Sparse cut $\Rightarrow$ small $\lambda_2$. Cheeger's inequality supplies the converse, which is the harder direction: small $\lambda_2$ \emph{forces} a sparse cut to exist, and you can find it by rounding the eigenvector.
\end{intuitionbox}

\begin{theorem}[Cheeger's inequality, discrete form]
\label{ch34:thm-cheeger}
For any connected $d$-regular graph $\Gph$,
\[
  \frac{\lambda_2}{2}\;\le\; h(\Gph)\;\le\;\sqrt{2\,d\,\lambda_2}.
\]
Equivalently, in conductance $\phi=h/d$: $\ \lambda_2/(2d)\le \phi\le \sqrt{2\lambda_2/d}$.
\end{theorem}

The two sides have very different characters. The left inequality (easy) says expansion can be no larger than $\lambda_2$ permits. The right inequality (the Cheeger rounding) says that if $\lambda_2$ is bounded below by a constant, then $h$ is too---which is exactly Definition~\ref{ch34:def-expander}. So:
\begin{center}
\emph{A family of $d$-regular graphs is a family of expanders iff $\lambda_2\ge \eps'$ for a constant $\eps'>0$.}
\end{center}
The spectral gap and combinatorial expansion are, up to the quadratic gap, the same phenomenon.

\begin{whybox}
\textbf{Proof of the easy direction $\lambda_2\le 2h$.} Let $S$ achieve $h=\abs{\partial S}/\abs S$ with $\abs S\le n/2$. Build the test vector
\[
  x_v=\begin{cases}\ \ \abs{\bar S} & v\in S,\\[2pt] -\abs S & v\in\bar S.\end{cases}
\]
Then $\sum_v x_v=\abs S\abs{\bar S}-\abs{\bar S}\abs S=0$, so $x\perp\mathbf 1$. An edge inside $S$ or inside $\bar S$ contributes $0$ to $x^\top\Lap x$; a boundary edge contributes $(\abs{\bar S}+\abs S)^2=n^2$. Hence $x^\top\Lap x=\abs{\partial S}\,n^2$. Meanwhile $x^\top x=\abs S\,\abs{\bar S}^2+\abs{\bar S}\,\abs S^2=\abs S\abs{\bar S}\,n$. By Courant--Fischer,
\[
  \lambda_2\le \frac{x^\top\Lap x}{x^\top x}=\frac{\abs{\partial S}\,n^2}{\abs S\abs{\bar S}\,n}=\frac{\abs{\partial S}\,n}{\abs S\abs{\bar S}}.
\]
Since $\abs{\bar S}\ge n/2$, we get $n/\abs{\bar S}\le 2$, so $\lambda_2\le 2\abs{\partial S}/\abs S=2h$. \hfill$\square$
\end{whybox}

The hard direction takes the second eigenvector $x$ (the Fiedler vector), sorts vertices by their coordinate value, and considers only the $n-1$ \emph{threshold cuts} $S_t=\set{v:x_v>t}$. A clever averaging argument---multiply and divide by $\sum\abs{x_u^2-x_v^2}$ and apply Cauchy--Schwarz across edges---shows at least one threshold cut has $\abs{\partial S_t}/\abs{S_t}\le\sqrt{2d\lambda_2}$. Algorithm~\ref{ch34:alg-cheeger} writes out the sweep procedure implicit in this proof, since it is what you actually run.

\begin{figure}[H]
\centering
\begin{tikzpicture}[scale=1.0,
  v/.style={circle,draw=nodeblue,fill=nodeblue!15,inner sep=1.4pt,font=\tiny}]
  % sorted Fiedler coordinates as a bar/line
  \def\xs{-1.7,-1.1,-0.6,-0.2,0.1,0.5,0.9,1.5}
  \foreach \x [count=\i from 1] in \xs {
    \node[v] (p\i) at ({(\i-4.5)*0.9},{\x}) {};
  }
  \draw[edgegray,-{Stealth}] (-3.9,-2.1) -- (3.9,-2.1) node[right,font=\scriptsize] {sorted vertices $v_1,\dots,v_n$};
  \draw[edgegray,-{Stealth}] (-3.9,-2.1) -- (-3.9,1.9) node[above,font=\scriptsize] {$x_{v}$};
  \foreach \i in {1,...,8}{\draw[edgegray!40,dashed] (p\i) -- ({(\i-4.5)*0.9},-2.1);}
  % one threshold cut between v4 and v5 (at x = 0.0)
  \draw[trapred,line width=1.2pt] (0.0,-2.25) -- (0.0,1.6);
  \node[trapred,font=\scriptsize] at (0.0,1.85) {threshold $t$};
  \node[font=\scriptsize] at (-2.0,1.3) {$S_t=\set{v:x_v\le t}$};
  \node[font=\scriptsize] at (1.5,1.3) {$\bar S_t=\set{v:x_v> t}$};
\end{tikzpicture}
\caption{Cheeger rounding (Algorithm~\ref{ch34:alg-cheeger}). Sort vertices by their Fiedler coordinate $x_v$; each of the $n-1$ gaps is a candidate threshold cut $S_t$. The sweep evaluates $\abs{\partial S_t}/\abs{S_t}$ at every gap and keeps the best; the analysis guarantees one of them is within the quadratic Cheeger bound $\sqrt{2d\lambda_2}$.}
\label{ch34:fig-sweep}
\end{figure}

\begin{algorithm}[H]
\caption{Spectral partitioning / Cheeger rounding}
\label{ch34:alg-cheeger}
\KwIn{connected $d$-regular graph $\Gph=(V,E)$ on $n$ vertices.}
\KwOut{a cut $S^\star$ with $\abs{\partial S^\star}/\abs{S^\star}\le\sqrt{2d\lambda_2}$.}
Form $\Lap=dI-\Adjm$\;
Compute $\lambda_2$ and its eigenvector $x$ (e.g.\ inverse power iteration deflating $\mathbf 1$)\;
Sort vertices so that $x_{v_1}\le x_{v_2}\le\cdots\le x_{v_n}$\;
$\mathit{best}\gets\infty$;\quad $S^\star\gets\emptyset$\;
\For{$t\gets 1$ \KwTo $n-1$}{
  $S_t\gets\set{v_1,\dots,v_t}$\;
  compute $\abs{\partial S_t}$ incrementally (add edges of $v_t$ to outside, subtract those now inside)\;
  $r\gets \abs{\partial S_t}/\min(\abs{S_t},n-\abs{S_t})$\;
  \If{$r<\mathit{best}$}{$\mathit{best}\gets r$;\quad $S^\star\gets S_t$\;}
}
\Return $S^\star$\;
\end{algorithm}

\begin{complexitybox}
The eigenvector dominates. With a sparse $\Lap$ ($dn$ nonzeros) and a near-linear-time Laplacian solver, $\lambda_2$ and $x$ come in $\tilde\Oh(dn)$ time; classical inverse iteration is $\Oh(\text{nnz}\cdot \text{iters})$. The sweep over threshold cuts (lines 5--11; see Figure~\ref{ch34:fig-sweep}) is $\Oh(dn)$ total, since each vertex is inserted once and each of its $d$ edges flips boundary status once. So the bottleneck is the eigenpair, and the guarantee on the returned cut is the quadratic Cheeger bound $\sqrt{2d\lambda_2}$.
\end{complexitybox}

\begin{trapbox}
The Cheeger guarantee is \emph{not} that you find the sparsest cut. You find a cut within a \emph{quadratic} factor of optimal: if the true $h$ is $\Th(\lambda_2)$ you are fine, but when $\lambda_2$ is tiny the rounding can be off by a $\sqrt{\,\cdot\,}$ factor. The quadratic loss is real and, in the worst case, unavoidable for this rounding---it is the reason Chapter~38 reaches for SDP-based sparsest-cut approximations when tighter guarantees matter. Spectral partitioning is fast and usually works well in practice, but it is not an exact cut oracle.
\end{trapbox}

\begin{marginnotebox}
The discrete Cheeger inequality, transported from Riemannian manifolds by Dodziuk and by Alon and Milman (1984--1985), makes $\lambda_2$ a computable certificate of expansion. Alon's 1986 paper ``Eigenvalues and expanders'' made this spectral connection a foundation of expander analysis and Markov chain mixing.
\end{marginnotebox}

\subsection{The adjacency spectrum and the second eigenvalue}
\label{ch34:adj-spectrum}

For $d$-regular graphs there is a parallel and equally common language using the \emph{adjacency} eigenvalues. Since $\Lap=dI-\Adjm$, if $\Adjm$ has eigenvalues $d=\mu_1\ge\mu_2\ge\cdots\ge\mu_n\ge -d$ then $\Lap$ has eigenvalues $\lambda_i=d-\mu_i$. The top adjacency eigenvalue is always $\mu_1=d$ (eigenvector $\mathbf 1$), corresponding to $\lambda_1=0$. The spectral gap $\lambda_2=d-\mu_2$ is large exactly when the second adjacency eigenvalue $\mu_2$ is much smaller than $d$.

Define
\[
  \mu(\Gph)\;=\;\max\bigl(\abs{\mu_2},\abs{\mu_n}\bigr),
\]
the largest \emph{nontrivial} adjacency eigenvalue in absolute value. Pulling in $\mu_n$ as well controls the bipartite-like ``near $-d$'' eigenvalues that ruin mixing. A $d$-regular graph is called a \emph{one-sided $(n,d,\lambda)$-expander} if $\mu_2\le\lambda$, and \emph{two-sided} if $\mu\le\lambda$ with $\lambda\ll d$. Smaller $\lambda$ means a better expander. The next two sections show why the ratio $\lambda/d$ is the parameter that matters.

\begin{formalbox}
\textbf{How small can $\lambda$ be?} Not arbitrarily. The \emph{Alon--Boppana bound} says any infinite family of $d$-regular graphs has $\mu_2\ge 2\sqrt{d-1}-o(1)$. So $2\sqrt{d-1}$ is a wall: no family beats it asymptotically. Graphs that meet it, $\mu=\max(\abs{\mu_2},\abs{\mu_n})\le 2\sqrt{d-1}$, are called \emph{Ramanujan graphs}---optimal expanders. They exist (Lubotzky--Phillips--Sarnak, 1988; Marcus--Spielman--Srivastava for bipartite, 2013), and they are as good as a sparse $d$-regular graph can ever be.
\end{formalbox}

Figure~\ref{ch34:fig-ramanujan-spectrum} shows this geometry on the adjacency eigenvalue axis $[-d,d]$: the trivial eigenvalue sits at $\mu_1=d$, while a two-sided expander traps all remaining eigenvalues $\mu_2,\dots,\mu_n$ inside $[-\mu,\mu]$ with a wide spectral gap $d-\mu$ on both ends.

\begin{figure}[H]
\centering
\begin{tikzpicture}[>=Stealth, font=\scriptsize]
  % Shaded Ramanujan bulk
  \fill[sourceblue!12,rounded corners=2pt] (-2.2,-0.32) rectangle (2.2,0.32);
  \draw[sourceblue,dashed,thick] (-2.2,-0.4) -- (-2.2,0.45);
  \draw[sourceblue,dashed,thick] (2.2,-0.4) -- (2.2,0.45);

  % Main axis
  \draw[edgegray,thick,->] (-4.8,0) -- (5.0,0) node[right,font=\small] {$\mu$};

  % Endpoints -d and +d
  \draw[trapred,very thick] (-4.2,-0.25) -- (-4.2,0.25);
  \node[below=4pt,trapred] at (-4.2,-0.25) {$-d$ (bipartite)};

  \draw[targetgreen,very thick] (4.2,-0.3) -- (4.2,0.4);
  \node[above=2pt,targetgreen,font=\scriptsize\bfseries] at (4.2,0.4) {$\mu_1 = d$};
  \node[below=4pt,targetgreen] at (4.2,-0.25) {$+d$ (trivial $\mathbf{1}$)};

  % Bulk ticks inside [-2\sqrt{d-1}, 2\sqrt{d-1}]
  \foreach \x in {-1.9,-1.4,-0.8,-0.3,0.2,0.7,1.3,1.85} {
    \draw[sourceblue,very thick] (\x,-0.2) -- (\x,0.2);
  }
  \node[above=2pt,sourceblue,font=\scriptsize\bfseries] at (0,0.42)
    {nontrivial bulk $\mu_2,\dots,\mu_n \in [-2\sqrt{d-1},\, 2\sqrt{d-1}]$};
  \node[below=4pt,sourceblue] at (-2.0,-0.35) {$-2\sqrt{d-1}$};
  \node[below=4pt,sourceblue] at (2.0,-0.35) {$+2\sqrt{d-1}$};

  % Spectral gap arrow
  \draw[<->,quotientorange,very thick] (2.25,0.18) -- (4.15,0.18);
  \node[above=2pt,quotientorange,font=\scriptsize\bfseries] at (3.2,0.2) {gap $d-\mu$};
\end{tikzpicture}
\caption{Adjacency spectrum of a $d$-regular Ramanujan expander. All nontrivial eigenvalues $\mu_2,\dots,\mu_n$ are confined within the Alon--Boppana interval $[-2\sqrt{d-1},2\sqrt{d-1}]$, keeping both the unmixing gap $d-\mu_2$ and the bipartite gap $d-\abs{\mu_n}$ large.}
\label{ch34:fig-ramanujan-spectrum}
\end{figure}

\section{The Expander Mixing Lemma}
\label{ch34:eml}

The central quantitative consequence of a spectral gap is that an expander graph \emph{mimics a random graph}. In a truly random $d$-regular graph, the number of edges between two sets $S,T$ should be about $\frac{d}{n}\abs S\abs T$ (each of the $\abs S\abs T$ ordered pairs is an edge with probability $\approx d/n$). The Expander Mixing Lemma says that in \emph{any} graph with small $\lambda$, the edge counts are close to that ideal, with error controlled by $\lambda$.

\begin{theorem}[Expander Mixing Lemma]
\label{ch34:thm-eml}
Let $\Gph$ be $d$-regular with $\mu=\max(\abs{\mu_2},\abs{\mu_n})$. For all $S,T\subseteq V$,
\[
  \abs[\Big]{\,e(S,T)-\frac{d}{n}\abs S\,\abs T\,}\;\le\;\mu\sqrt{\abs S\,\abs T},
\]
where $e(S,T)$ counts edges with one endpoint in $S$ and one in $T$ (pairs in $S\cap T$ counted twice).
\end{theorem}

\begin{whybox}
\textbf{Proof.} Let $\mathbf 1_S,\mathbf 1_T$ be indicator vectors. Then $e(S,T)=\mathbf 1_S^\top\Adjm\,\mathbf 1_T$ exactly, by definition of $\Adjm$. Decompose each indicator into its component along $\mathbf 1$ and the rest:
\[
  \mathbf 1_S=\frac{\abs S}{n}\mathbf 1 + s^\perp,\qquad \mathbf 1_T=\frac{\abs T}{n}\mathbf 1 + t^\perp,
\]
with $s^\perp,t^\perp\perp\mathbf 1$. Since $\Adjm\mathbf 1=d\mathbf 1$,
\[
  \mathbf 1_S^\top\Adjm\,\mathbf 1_T
  =\frac{\abs S\abs T}{n^2}\,\mathbf 1^\top\Adjm\mathbf 1 + (s^\perp)^\top\Adjm\,t^\perp
  =\frac{d}{n}\abs S\abs T + (s^\perp)^\top\Adjm\,t^\perp .
\]
The cross terms vanish because $\Adjm\mathbf 1=d\mathbf 1\perp s^\perp,t^\perp$. On the subspace orthogonal to $\mathbf 1$, every eigenvalue of $\Adjm$ has absolute value $\le\mu$, so $\norm{\Adjm\,t^\perp}\le\mu\norm{t^\perp}$. By Cauchy--Schwarz,
\[
  \abs{(s^\perp)^\top\Adjm\,t^\perp}\le\mu\,\norm{s^\perp}\norm{t^\perp}\le\mu\sqrt{\abs S}\sqrt{\abs T},
\]
using $\norm{s^\perp}^2=\abs S(1-\abs S/n)\le\abs S$ and likewise for $T$. Rearranging gives the claim. \hfill$\square$
\end{whybox}

The proof is three lines once you see that $e(S,T)$ is literally a quadratic form in the adjacency matrix and that orthogonality to $\mathbf 1$ kills the leading term. This is the cleanest illustration of the chapter's thesis: \emph{combinatorial counting becomes linear algebra}.

\begin{labbox}
Test the lemma on the cube $Q_3$ ($d=3$, $n=8$). Its adjacency eigenvalues are $3,1,1,1,-1,-1,-1,-3$, so $\mu=\max(\abs 1,\abs{-3})=3$. That equals $d$, so $Q_3$ is a \emph{bad} expander (the $-3$ eigenvalue flags that it is bipartite). Indeed for $S=$ even-weight strings, $T=$ odd-weight strings ($\abs S=\abs T=4$), every edge of $Q_3$ goes between them: $e(S,T)=12$, while $\frac{d}{n}\abs S\abs T=\frac{3}{8}\cdot 16=6$. The deviation is $6$, against the stated bound $\mu\sqrt{\abs S\abs T}=3\cdot 4=12$---so the lemma holds with a factor-of-two cushion. The slack is exactly the relaxation $\norm{s^\perp}^2\le\abs S$ in the proof: here $\norm{s^\perp}^2=\abs S(1-\abs S/n)=4\cdot\tfrac12=2$, and the tighter Cauchy--Schwarz form $\mu\norm{s^\perp}\norm{t^\perp}=3\cdot\sqrt2\cdot\sqrt2=6$ \emph{is} met with equality, driven by the $\mu_n=-3$ eigenvector (the bipartition sign vector). The lesson: bipartiteness ($\mu_n=-d$) obstructs rapid mixing, which is why the two-sided spectral bound depends on the absolute value $\abs{\mu_n}$ alongside $\mu_2$.
\end{labbox}

\begin{corollary}[Spectral expansion controls edge expansion]
\label{ch34:cor-eml-cut}
For $\abs S\le n/2$, setting $T=\bar S$ in Theorem~\ref{ch34:thm-eml} gives
\[
  \abs{\partial S}=e(S,\bar S)\ge \frac{d}{n}\abs S\abs{\bar S}-\mu\sqrt{\abs S\abs{\bar S}}.
\]
When $\mu\ll d$ this already pins down balanced cuts: at $\abs S=\abs{\bar S}=n/2$ it reads $\abs{\partial S}\ge \bigl(\tfrac{d}{2}-\mu\bigr)\tfrac n2$. For the uniform statement over \emph{all} $\abs S\le n/2$ it is cleaner to route through the easy Cheeger direction: $d-\mu_2=\lambda_2\le 2h$, and since $\mu_2\le\mu$,
\[
  h(\Gph)\ \ge\ \frac{\lambda_2}{2}\ =\ \frac{d-\mu_2}{2}\ \ge\ \frac{d-\mu}{2}.
\]
A spectral gap $d-\mu$ bounded away from $0$ forces constant edge expansion.
\end{corollary}

\section{Random graphs are expanders}
\label{ch34:random}

Existence is easy and instructive: almost every sparse regular graph is a near-optimal expander. The cleanest model is the \emph{random $d$-regular graph} $\Gph(n,d)$, drawn uniformly from all $d$-regular graphs on $n$ labeled vertices (concretely, via the configuration / pairing model: give each vertex $d$ half-edges and match the $dn$ half-edges uniformly at random). We prove combinatorial expansion directly with a union bound; the spectral statement is deeper.

\begin{theorem}[A random regular graph expands]
\label{ch34:thm-random-expand}
Fix $d\ge 3$. There is a constant $\eps>0$ (depending only on $d$) such that a uniformly random $d$-regular graph on $n$ vertices satisfies $h(\Gph)\ge\eps$ with probability $1-o(1)$.
\end{theorem}

\begin{whybox}
\textbf{Proof sketch (first-moment / union bound).} Call a set $S$ with $\abs S=k\le n/2$ \emph{bad} if $\abs{\partial S}<\eps k$, i.e.\ fewer than $\eps k$ edges leave it. We bound the expected number of bad sets and show it is $o(1)$ when $\eps$ is small enough.

Fix $k$. There are $\binom{n}{k}$ candidate sets. For a fixed $S$ of size $k$ to be bad in the pairing model, most of its $dk$ half-edges must match \emph{within} $S$: at most $\eps k$ may cross to $\bar S$. The probability that a given half-edge of $S$ stays inside is roughly $dk/(dn)=k/n$, and the events are nearly independent, so the chance that all but $\eps k$ of the $dk$ half-edges stay inside is at most
\[
  \binom{dk}{\eps k}\left(\frac{k}{n}\right)^{dk-\eps k}.
\]
Multiply by $\binom nk\le (en/k)^k$ and sum over $k$. The dominant factor is $(en/k)^k\cdot (k/n)^{(d-\eps)k}=\bigl[e\,(k/n)^{\,d-\eps-1}\bigr]^{k}$. For $d\ge 3$ and $\eps$ small the exponent $d-\eps-1\ge 1$, so $(k/n)^{d-\eps-1}\le k/n$, and the bracket is at most $e\,k/n<1$ for $k\le n/2$ once we absorb the binomial entropy. The geometric-type sum over $k$ is $o(1)$. Hence with probability $1-o(1)$ no bad set exists, i.e.\ $h\ge\eps$. \hfill$\square$
\end{whybox}

\begin{intuitionbox}
Why does randomness help? A small set $S$ has $dk$ half-edges waiting to be paired. There are $dn$ half-edges total but only $dk$ inside $S$. With $k\le n/2$, the overwhelming majority of partners live outside $S$, so a typical half-edge escapes. To trap a set (make it bad) you would need a wildly atypical matching that keeps almost all of $S$'s half-edges internal---combinatorially rare. Expansion is the \emph{generic} behavior; non-expanders (cycles, grids) are the special, structured exceptions.
\end{intuitionbox}

The spectral counterpart is the celebrated result of Friedman (2008), settling Alon's conjecture: a random $d$-regular graph is \emph{almost Ramanujan}.

\begin{theorem}[Friedman, informal]
\label{ch34:thm-friedman}
For every $\eps>0$, a random $d$-regular graph satisfies $\mu(\Gph)\le 2\sqrt{d-1}+\eps$ with probability $1-o(1)$.
\end{theorem}

So random graphs essentially meet the Alon--Boppana wall. This settles existence but does not help when you need a \emph{specific}, checkable graph---e.g.\ wired into hardware or a code---and cannot afford to verify $\lambda_2$ on the fly or risk a bad draw. That is the motivation for explicit constructions.

\begin{table}[H]
\centering
\begin{tabular}{lcccc}
\toprule
 & degree $d$ & edges & $h(\Gph)$ as $n\to\infty$ & conductance $\phi=h/d$ \\
\midrule
cycle $C_n$        & $2$ & $n$        & $\to 0\ (\sim 4/n)$ & $\to 0\ (\sim 2/n)$ \\
$\sqrt n\times\sqrt n$ grid & $4$ & $\sim 2n$ & $\to 0\ (\sim 1/\sqrt n)$ & $\to 0$ \\
hypercube $Q_{\log n}$ & $\log n$ & $\frac{n\log n}{2}$ & $1$ (half-cube cut) & $\Th(1/\log n)$ \\
random $d$-regular & $d$ (fixed) & $dn/2$ & $\Th(1)$ & $\Th(1)$ \\
complete $K_n$     & $n-1$ & $\binom n2$ & $\Th(n)$ & $\to 1/2$ \\
\bottomrule
\end{tabular}
\caption{Expansion across familiar graphs. Edge expansion $h$ is not degree-normalized, so the degree-fair comparison is the conductance $\phi=h/d\in[0,1]$. Only the random regular graph keeps $\phi$ bounded away from zero \emph{at fixed degree and linear edge count}---the expander sweet spot. The hypercube has constant $h$ but only because its degree $\log n$ grows, and its conductance still decays.}
\label{ch34:tab-zoo}
\end{table}

\section{Explicit constructions, conceptually}
\label{ch34:constructions}

A probabilistic existence proof gives an object you cannot point to. For many uses---hardware interconnects, deterministic algorithms, codes you ship---you need an \emph{explicit} family: given $n$, write down the graph (or even just answer ``who are vertex $v$'s neighbors?'') in $\poly(\log n)$ time, and \emph{prove} its expansion once and for all. There are three main approaches.

\paragraph{Algebraic constructions.} The earliest explicit expanders come from group theory and number theory. A clean example is the \emph{Margulis} construction: take vertices $\Z_m\times\Z_m$ and connect each $(x,y)$ to a fixed handful of affine images such as $(x\pm y,y),(x,y\pm x),(x\pm y+1,y),\dots$. The eigenvalue bound is proved using harmonic analysis on the group. The strongest algebraic examples are the \emph{Lubotzky--Phillips--Sarnak} (LPS) graphs: Cayley graphs of $\mathrm{PGL}_2$ over finite fields whose Ramanujan bound $\mu\le 2\sqrt{d-1}$ follows from deep number theory (Deligne's proof of the Ramanujan--Petersson conjecture). These are optimal but the machinery is heavy.

\paragraph{The zig-zag product.} Reingold, Vadhan, and Wigderson (2002) gave a purely combinatorial, self-contained route. The \emph{zig-zag product} $\Gph\mathbin{\text{\textcircled z}} H$ combines a big graph $\Gph$ (good expansion, large degree) with a small graph $H$ (fixed size, good expansion) to produce a graph on $\abs{V(\Gph)}\cdot\abs{V(H)}$ vertices whose degree is that of $H$ (small!) and whose spectral gap degrades only mildly.

\begin{intuitionbox}
Replace every vertex of $\Gph$ by a copy (``cloud'') of $H$, one cloud-vertex per edge of $\Gph$. A step in the product is \emph{zig--zag--zig}: take a small step inside your cloud (using $H$), then traverse the $\Gph$-edge you landed on to jump to another cloud, then a small step inside the new cloud. The two little $H$-steps ``randomize'' which $\Gph$-edge you cross, so even though each vertex now has only $\deg(H)^2$ neighbors, the long-range mixing of $\Gph$ is inherited. You trade $\Gph$'s big degree for $H$'s small degree while keeping expansion. Iterating (square $\Gph$ to boost the gap before each zig-zag) builds an infinite family from one small gadget $H$.
\end{intuitionbox}

\begin{figure}[H]
\centering
\begin{tikzpicture}[scale=1.0,
  cl/.style={circle,draw=structurepurple,dashed,inner sep=10pt},
  v/.style={circle,draw=nodeblue,fill=nodeblue!15,inner sep=1.3pt}]
  \node[cl] (A) at (0,0) {};
  \node[cl] (B) at (4,0) {};
  \node at (0,-1.2) {\scriptsize cloud (copy of $H$)};
  \node at (4,-1.2) {\scriptsize cloud (copy of $H$)};
  \foreach \a in {0,1,2}{\node[v] (a\a) at ({120*\a+90}:0.55) {};}
  \begin{scope}[shift={(4,0)}]
    \foreach \b in {0,1,2}{\node[v] (b\b) at ({120*\b+90}:0.55) {};}
  \end{scope}
  \draw[edgegray!50] (a0)--(a1)--(a2)--(a0);
  \draw[edgegray!50] (b0)--(b1)--(b2)--(b0);
  \draw[targetgreen,line width=1.2pt] (a1)--(b2) node[midway,above,font=\scriptsize] {step 2: $\Gph$-edge};
  \draw[sourceblue,-{Stealth},line width=1pt] (a0) to[bend left=30] (a1);
  \node[sourceblue,font=\scriptsize,anchor=east] at (-0.75,0.55) {step 1: zig in $H$};
  \draw[trapred,-{Stealth},line width=1pt] (b2) to[bend left=30] (b0);
  \node[trapred,font=\scriptsize,anchor=west] at (4.75,0.55) {step 3: zag in $H$};
\end{tikzpicture}
\caption{One step of the zig-zag product: a small $H$-step (zig), a $\Gph$-edge crossing between clouds, then another $H$-step (zag). The two $H$-steps randomize the crossing edge, so a tiny degree suffices.}
\label{ch34:fig-zigzag}
\end{figure}

\paragraph{Lossless / bipartite expanders.} For coding theory one wants \emph{vertex} expansion close to the degree: every small left-set $S$ has $\abs{N(S)}\ge(1-\eps)d\abs S$ neighbors, nearly as many as the $d\abs S$ edges leaving it. These ``lossless'' expanders (Capalbo--Reingold--Vadhan--Wigderson, 2002) are built from a randomness-conductor variant of zig-zag and are exactly what Section~\ref{ch34:codes} needs. Spectral methods alone cannot reach $A$ close to $d$ (the Kahale barrier), so these need the combinatorial constructions.

\begin{marginnotebox}
Explicit expanders evolved in stages: Margulis gave the first algebraic construction in 1973, Gabber and Galil made bounds explicit in 1981, and Lubotzky--Phillips--Sarnak reached spectral optimality in 1988. The zig-zag product (2002) introduced a purely combinatorial path, leading directly to Reingold's proof that $\mathsf{SL}=\mathsf{L}$ (undirected reachability in deterministic logspace). In 2013, Marcus, Spielman, and Srivastava proved the existence of bipartite Ramanujan graphs of every degree via interlacing polynomials.
\end{marginnotebox}

\section{Why expanders matter}
\label{ch34:applications}

We now turn to four applications, each using a different consequence of the spectral gap.

\subsection{Fast mixing of random walks}
\label{ch34:mixing}

Run a lazy random walk on a $d$-regular graph: from $v$, stay put with probability $1/2$, else step to a uniform neighbor. The transition matrix is $W=\tfrac12 I+\tfrac12\tfrac1d\Adjm$, with eigenvalues $w_i=\tfrac12+\tfrac{\mu_i}{2d}\in[0,1]$. The stationary distribution is uniform, $\pi=\mathbf 1/n$. How fast does the walk forget its start?

\begin{theorem}[Mixing time of an expander]
\label{ch34:thm-mixing}
Let $w_\star=\max_{i\ge2}\abs{w_i}$ be the largest nontrivial walk eigenvalue in absolute value. Laziness keeps every $w_i=\tfrac12+\tfrac{\mu_i}{2d}\in[0,1]$, so the maximum is attained at the most positive nontrivial adjacency eigenvalue $\mu_2$:
\[
  w_\star=\tfrac12+\tfrac{\mu_2}{2d}\ \le\ \tfrac12+\tfrac{\mu}{2d}\ <\ 1 .
\]
For any start $v$ and any $t$, the distribution $p_t$ satisfies
\[
  \norm{p_t-\pi}_1\ \le\ \sqrt n\,\cdot\, w_\star^{\,t}.
\]
Hence after $t=\Oh\!\bigl(\tfrac{1}{1-w_\star}\log n\bigr)=\Oh\!\bigl(\tfrac{d}{d-\mu_2}\log n\bigr)$ steps the walk is within $1/n$ of uniform in total variation. (Replacing $\mu_2$ by $\mu$ only loosens the bound, since $\mu_2\le\mu$.)
\end{theorem}

\begin{whybox}
\textbf{Why.} Write the start vector $e_v=\pi+\sum_{i\ge2}c_i\psi_i$ in the eigenbasis $\set{\psi_i}$ of $W$. Applying $W$ once scales the $i$-th component by $w_i$; the $\pi$-component (eigenvalue $1$) is fixed. So
\[
  \norm{p_t-\pi}_2^2=\sum_{i\ge2}c_i^2 w_i^{2t}\le w_\star^{2t}\sum_{i\ge2}c_i^2\le w_\star^{2t}.
\]
Convert $\ell_2$ to $\ell_1$ with $\norm{\,\cdot\,}_1\le\sqrt n\,\norm{\,\cdot\,}_2$. The decay rate is governed entirely by the \emph{spectral gap} $1-w_\star=\tfrac{d-\mu_2}{2d}$: a constant gap (expander!) gives $\Oh(\log n)$ mixing. \hfill$\square$
\end{whybox}

\begin{intuitionbox}
On a cycle the walk diffuses like Brownian motion, reaching the far side only after $\Th(n^2)$ steps---the gap is $\Th(1/n^2)$. On an expander every region leaks into every other immediately, so $\Th(\log n)$ steps suffice. Fast mixing is the dynamical counterpart of expansion, and for sampling, $\log n$ versus $n^2$ steps is the difference between feasible and infeasible.
\end{intuitionbox}

\begin{applicationbox}
Fast mixing is what makes \emph{Markov chain Monte Carlo} work (Chapters~35 and~36). To sample from a complicated distribution---a random matching, a random coloring, a Gibbs state---you design a chain whose stationary distribution is the target and whose state graph is an expander; rapid mixing then means a short, provably good sample. The same gap bounds card-shuffling times and the convergence of PageRank.
\end{applicationbox}

\subsection{Robust networks}
\label{ch34:robust}

Because every cut is large, expanders are hard to disconnect and hard to congest. Edge expansion $h\ge\eps$ means you must delete at least $\eps$ edges \emph{per vertex} to peel off any region: the edge-connectivity is high relative to the sparse degree, so random failures and even adversarial attacks leave a giant connected core. Expanders also make good \emph{routing} networks: any permutation of $n$ packets can be routed with low congestion and dilation $\Oh(\log n)$, because the mixing property spreads load uniformly. This is why expander topologies appear in data-center and supercomputer interconnects and in peer-to-peer overlays, where you want maximal resilience and bandwidth from a bounded number of links per node.

\subsection{Derandomization}
\label{ch34:derand}

Suppose a randomized algorithm decides correctly with probability $2/3$ using $r$ random bits, and you want error $2^{-k}$. Naively you repeat $\Oh(k)$ times, spending $\Oh(kr)$ fresh random bits. Random bits are an expensive resource in complexity theory; can we spend fewer?

\begin{intuitionbox}
Build the \emph{expander walk}. Make every $r$-bit string a vertex of a constant-degree expander on $N=2^r$ vertices. Pick a start uniformly ($r$ bits), then take a $k$-step random walk, each step costing only $\log d=\Oh(1)$ bits to choose a neighbor. Total randomness: $r+\Oh(k)$ bits, not $\Oh(kr)$. The walk visits $k$ correlated-but-spread-out strings; use them as the $k$ random seeds.
\end{intuitionbox}

\begin{theorem}[Ajtai--Koml\'os--Szemer\'edi / Gillman, expander Chernoff bound]
\label{ch34:thm-chernoff}
Let $B\subseteq V$ be a ``bad'' set (seeds where the algorithm errs) with $\abs B/n\le 1/3$, in a graph with $\mu/d=\beta$. For a random walk $v_1,\dots,v_k$,
\[
  \Pr\!\Bigl[\#\set{i:v_i\in B}\ge \tfrac12 k\Bigr]\ \le\ e^{-\Om((1-\beta)k)} .
\]
\end{theorem}

So the fraction of walk steps landing in the bad set concentrates just as if the steps were independent, with a rate set by the spectral gap $1-\beta$. Taking the majority vote of the $k$ runs errs with probability $2^{-\Om(k)}$---the same amplification as $k$ independent trials, at a fraction of the randomness. This expander-walk sampler is a standard tool in derandomization and underlies the proof that error reduction costs essentially no extra random bits.

\subsection{Error-correcting codes}
\label{ch34:codes}

Expanders give codes that encode, and especially \emph{decode}, in linear time. Sipser and Spielman (1996) defined \emph{expander codes} from a bipartite expander: $n$ variable nodes on the left, $m$ constraint nodes on the right, each constraint forcing the parity (or a small sub-code) of its neighborhood to be valid. A codeword is an assignment satisfying all constraints.

\begin{whybox}
\textbf{Why expansion decodes.} Use a lossless left expander: every small set $S$ of corrupted variables has $\abs{N(S)}\ge(1-\eps)d\abs S$ neighbors with $\eps<1/4$. Then most constraints adjacent to $S$ touch $S$ \emph{only once}, so they are \emph{unsatisfied} and visibly flag the error. The decoder is greedy: \emph{find any variable that would fix more constraints than it breaks, and flip it.} Each flip strictly decreases the number of unsatisfied constraints, so the process terminates in linear time, and expansion guarantees that as long as the error set is small such a beneficial flip always exists until all errors are gone.
\end{whybox}

\begin{theorem}[Sipser--Spielman, informal]
\label{ch34:thm-expander-codes}
A code from a $(\gamma n,\,(1-\eps)d)$ lossless bipartite expander with $\eps<1/4$ has constant rate, corrects a constant fraction of adversarial errors, and admits a flip-decoder running in $\Oh(n)$ time.
\end{theorem}

\begin{figure}[H]
\centering
\begin{tikzpicture}[scale=1,
  var/.style={circle,draw=nodeblue,fill=nodeblue!15,inner sep=2pt},
  con/.style={rectangle,draw=quotientorange,fill=quotientorange!15,inner sep=3pt}]
  \foreach \i in {1,...,5}{\node[var] (x\i) at (0,{4-0.8*\i}) {};}
  \node at (0,3.7) {\scriptsize variables};
  \foreach \j in {1,...,3}{\node[con] (c\j) at (3,{3.4-1.0*\j}) {$+$};}
  \node at (3,3.7) {\scriptsize parity checks};
  \draw (x1)--(c1) (x2)--(c1) (x3)--(c1);
  \draw (x2)--(c2) (x3)--(c2) (x4)--(c2);
  \draw (x3)--(c3) (x4)--(c3) (x5)--(c3);
  \node[var,fill=trapred!30,draw=trapred] at (0,3.2) {};
  \node[trapred,font=\scriptsize] at (-1.4,3.2) {flipped bit};
\end{tikzpicture}
\caption{An expander (Tanner) code. Left vertices are bits, right vertices parity checks. A corrupted bit makes several checks fail; lossless expansion guarantees a greedy bit-flip always exists until decoding succeeds.}
\label{ch34:fig-code}
\end{figure}

\begin{connectionbox}
Each part of this chapter connects back to Chapter~31 and the companion \emph{Linear Algebra} volume. Expansion $=$ spectral gap $=$ fast mixing $=$ small effective resistances is one phenomenon viewed four ways. The Laplacian quadratic form (Chapter~31) gives Cheeger; the adjacency spectrum gives the Mixing Lemma; the walk matrix $W=\tfrac12(I+\tfrac1d\Adjm)$ gives mixing time; and in the companion \emph{Linear Algebra} volume the same gap appears as the condition number of the normalized Laplacian, which is precisely what makes expander \emph{sparsifiers} and electrical-flow solvers fast. For expanders, the combinatorial, probabilistic, and spectral descriptions all express one property.
\end{connectionbox}

\section{A fully worked example}
\label{ch34:worked}

Let us run the whole pipeline on one concrete graph: the \emph{Petersen graph} $P$, which is $3$-regular on $10$ vertices and a standard small expander.

\begin{intuitionbox}
$P$ has an outer pentagon $0,1,2,3,4$, an inner pentagram $5,6,7,8,9$, and spokes $i$--$(i+5)$. It is vertex-transitive, has girth $5$, and is a very good small expander. Its adjacency eigenvalues are $3$ (once), $1$ (five times), and $-2$ (four times). So $\mu=\max(\abs 1,\abs{-2})=2$, the spectral gap is $d-\mu=1$, and $\lambda_2=d-\mu_2=3-1=2$.
\end{intuitionbox}

\begin{figure}[H]
\centering
\begin{tikzpicture}[scale=1.5,
  v/.style={circle,draw=nodeblue,fill=nodeblue!15,inner sep=1.5pt,font=\tiny}]
  \foreach \i in {0,...,4}{\node[v] (o\i) at ({90+72*\i}:1) {\i};}
  \foreach \i in {0,...,4}{\node[v] (n\i) at ({90+72*\i}:0.5) {\the\numexpr\i+5\relax};}
  \foreach \i [evaluate=\i as \j using {int(mod(\i+1,5))}] in {0,...,4}{\draw (o\i)--(o\j);}
  \foreach \i [evaluate=\i as \j using {int(mod(\i+2,5))}] in {0,...,4}{\draw (n\i)--(n\j);}
  \foreach \i in {0,...,4}{\draw (o\i)--(n\i);}
\end{tikzpicture}
\caption{The Petersen graph: outer $5$-cycle, inner pentagram, five spokes. A small $(10,3,2)$-expander, in fact Ramanujan since $\mu=2\le 2\sqrt{3-1}=2\sqrt2\approx 2.83$.}
\label{ch34:fig-petersen}
\end{figure}

\textbf{Step 1: edge expansion by hand.} By vertex-transitivity it suffices to test representative cuts. A single vertex has $\abs{\partial S}/\abs S=3$. A spoke pair like $\set{0,5}$: vertex $0$ has neighbors $1,4,5$ and vertex $5$ has neighbors $0,7,8$; the boundary is $\set{1,4,7,8}$ via four edges (edge $0$--$5$ is internal), ratio $4/2=2$. An outer edge $\set{0,1}$: boundaries to $4,5$ (from $0$) and $2,6$ (from $1$), four edges, ratio $2$. The sparsest balanced cuts of $P$ achieve $h(P)=1$: a $5$/$5$ split can leave only $5$ edges crossing, giving ratio $5/5=1$, and one checks the few balanced cuts that none drops below that.

\textbf{Step 2: check Cheeger.} Theorem~\ref{ch34:thm-cheeger} predicts
\[
  \frac{\lambda_2}{2}=\frac{2}{2}=1\ \le\ h(P)=1\ \le\ \sqrt{2d\lambda_2}=\sqrt{2\cdot 3\cdot 2}=\sqrt{12}\approx 3.46.
\]
Both inequalities hold, and the \emph{lower} bound is tight here: $h(P)=\lambda_2/2$ exactly. The upper Cheeger bound is loose by a factor $\approx 3.46$, the usual quadratic slack.

\textbf{Step 3: Corollary~\ref{ch34:cor-eml-cut}.} The spectral cut bound gives $h\ge(d-\mu)/2=(3-2)/2=1/2$, weaker than $1$ but in the right ballpark. Here $\mu=2$ is large relative to $d=3$, so the raw mixing-lemma count is slack for small sets; the clean $h\ge(d-\mu)/2$ comes through the easy Cheeger step $d-\mu_2=\lambda_2\le 2h$ (with $\mu_2=1\le\mu$). Using $\mu_2$ instead recovers the tight value $h\ge(d-\mu_2)/2=1$.

\textbf{Step 4: mixing.} The walk eigenvalues are $w_i=\tfrac12+\tfrac{\mu_i}{6}$, so $\mu_i=1$ gives $\tfrac23$ and $\mu_i=-2$ gives $\tfrac16$. The largest nontrivial one is $w_\star=\tfrac12+\tfrac{\mu_2}{2d}=\tfrac12+\tfrac16=\tfrac23$ (the negative eigenvalue is tamed by laziness into the small $\tfrac16$, not a near-$1$ value). The gap is $1-w_\star=1/3$, so from any start the walk is within $\Th(1/n)$ of uniform after $\Oh\!\bigl(\tfrac{d}{d-\mu_2}\log 10\bigr)=\Oh(3\log 10)\approx \Oh(7)$ steps---a tiny graph mixes almost instantly, as expansion predicts.

\begin{labbox}
Reproduce Step 4 numerically. Start $p_0=e_0$ (all mass on vertex $0$). One lazy step: $p_1$ has $1/2$ on vertex $0$ and $1/6$ on each of $1,4,5$. Iterate $p_{t+1}=Wp_t$ a dozen times (a $10\times10$ matrix multiply); watch $\norm{p_t-\tfrac1{10}\mathbf 1}_1$ fall by a factor near $2/3$ each step. Then deliberately break expansion: delete the five spokes, splitting $P$ into two disjoint $5$-cycles. Now $\mu=d$, the gap vanishes, and the walk never mixes across the (now nonexistent) cut. The experiment shows mixing failing as soon as the spectral gap closes.
\end{labbox}

\textbf{Step 5: watch the error fall.} By vertex-transitivity the only thing that distinguishes coordinates of $p_t$ is the distance from the start, so we can collapse the $10$-vector into three numbers: the mass on vertex $0$ (distance $0$), the mass on each of its three neighbors (distance $1$), and the mass on each of the six remaining vertices (distance $2$). The lazy step $p_{t+1}=Wp_t$ becomes a $3\times3$ recurrence, and the total-variation distance to uniform $\tfrac12\norm{p_t-\tfrac1{10}\mathbf 1}_1$ is now trivial to track by hand.

\begin{table}[H]
\centering
\begin{tabular}{rcccc}
\toprule
$t$ & mass at dist.\ $0$ & each at dist.\ $1$ & each at dist.\ $2$ & $\tfrac12\norm{p_t-\pi}_1$ \\
\midrule
$0$ & $1.000$ & $0.000$ & $0.000$ & $0.900$ \\
$1$ & $0.500$ & $0.167$ & $0.000$ & $0.600$ \\
$2$ & $0.333$ & $0.167$ & $0.028$ & $0.433$ \\
$3$ & $0.250$ & $0.148$ & $0.051$ & $0.294$ \\
$5$ & $0.166$ & $0.122$ & $0.078$ & $0.132$ \\
$8$ & $0.120$ & $0.107$ & $0.093$ & $0.039$ \\
$12$& $0.104$ & $0.101$ & $0.099$ & $0.008$ \\
\bottomrule
\end{tabular}
\caption{Lazy walk on the Petersen graph from vertex $0$, collapsed by distance. The total-variation distance to uniform contracts by roughly the predicted factor $w_\star=2/3$ per step once the transient settles, dropping below $1/n=0.1$ just after step $5$ (it is $0.088$ at step $6$) and reaching $0.008$ by step $12$.}
\label{ch34:tab-mixing}
\end{table}

Read the last column down: after the first couple of steps each entry is about $2/3$ of the one above it, exactly the eigenvalue $w_\star$ predicted in Step~4. The distance crosses the $1/n$ threshold just after $t=5$, comfortably inside the budget: the bound $\sqrt n\,w_\star^{\,t}\le 1/n$ asks for only about $9$ steps. Nothing in this table required diagonalizing a matrix; the spectral gap told us the contraction rate in advance, and the arithmetic merely confirmed it.

\section{Historical Context and Controversies}
\label{ch34:history-context}

\begin{historybox}
The story of expander graphs began in telecommunications. In $1973$, Mark Pinsker at the Institute for Information Transmission Problems in Moscow was studying telephone switching networks (routing systems that connect telephone callers without colliding). A telephone exchange needed to connect thousands of callers simultaneously without blocking, using as few internal switches and wires as possible. Pinsker realized that the ideal switching network was a sparse bipartite graph where every subset of inputs fanned out to a strictly larger set of outputs. Using the probabilistic method (proving an object exists by showing that a random choice succeeds with positive probability), Pinsker proved that a randomly chosen sparse graph is almost certainly an excellent expander.

However, Pinsker's proof was non-constructive: it showed that almost all sparse graphs are expanders, but gave no recipe to build even a single one deterministically. Wiring a physical exchange required an explicit blueprint. That same year, Gregory Margulis ($1973$) published the first explicit construction of constant-degree expanders. Because purely combinatorial techniques could not certify expansion, Margulis turned to representation theory, using Kazhdan's property~(T) (a rigid representation-theoretic property of groups where the trivial representation is isolated) to prove a positive spectral gap. In $1988$, Alexander Lubotzky, Ralph Phillips, and Peter Sarnak (LPS), and independently Margulis, constructed explicit \emph{Ramanujan graphs} (optimal expanders matching the Alon--Boppana optimal spectral bound $2\sqrt{d-1}$) using Pierre Deligne's proof of the Ramanujan conjecture in automorphic forms.

For decades, the search for explicit expanders posed a classic paradox: a haystack made almost entirely of needles, yet hard to exhibit a single needle deterministically without heavy algebraic machinery. While dense graphs like $K_n$ expand maximally, they overwhelm wiring and memory budgets with $O(n^2)$ edges; the engineering goal was to achieve constant degree $d = O(1)$ alongside linear boundary expansion. In $2002$, Omer Reingold, Salil Vadhan, and Avi Wigderson developed the \emph{zig-zag product} (an operation that replaces each vertex of a large graph with an entire small graph to amplify expansion without inflating degree), enabling the construction of constant-degree expanders via elementary graph-theoretic operations and leading directly to Reingold's $2005$ deterministic logspace connectivity algorithm ($\mathbf{L} = \mathbf{SL}$). An open problem is to construct explicit non-bipartite Ramanujan graphs for every degree and size using purely combinatorial methods, bypassing heavy algebraic number theory.
\end{historybox}

\subsection{What We Can Now Do and Its Limitations}

\begin{focusbox}
\textbf{Capabilities Unlocked}:
\begin{itemize}
  \item \textbf{Derandomization and Randomness Reduction}: Expanders allow randomized algorithms to be run with drastically fewer
        random bits. Rather than drawing $k$ independent random seeds of length $\log n$ (costing $k\log n$ bits), a random walk on an
        expander generates $k$ pseudorandom samples using only $\log n + \Oh(k)$ total bits while achieving nearly identical error decay.
  \item \textbf{Linear-Time Error-Correcting Codes}: Sipser and Spielman ($1996$) used bipartite expanders to construct error-correcting
        codes that can be encoded and decoded in strictly linear time.
  \item \textbf{Fast Information Diffusion}: Gossip protocols and routing networks built on expanders guarantee that information reaches
        every node in $O(\log n)$ rounds with no single point of failure.
\end{itemize}

\textbf{Honest Limitations and Downsides}:
\begin{itemize}
  \item \textbf{Rigid Algebraic Sizes}: Optimal Ramanujan constructions exist only for specific vertex counts and degrees related to
        primes (e.g., $d = p+1$ for prime $p \equiv 1 \pmod 4$), making them cumbersome for arbitrary system sizes.
  \item \textbf{Absence of Efficient Routing}: An expander guarantees that short paths exist between every pair of vertices, but it does
        \emph{not} provide a localized routing algorithm for packets to find those short paths without global topology tables.
\end{itemize}

\textbf{Where We Go Next in the Book}:
\begin{itemize}
  \item In \textbf{Chapter~35 (Random Walks and Electrical Networks)}, we connect the expander spectral gap directly to electrical
        currents, effective resistances, and Ohm's law.
  \item In \textbf{Chapter~36 (PageRank and Markov Chains)}, we use the spectral gap of the Google matrix to prove exponential convergence
        of power iteration on web graphs.
  \item In \textbf{Chapter~41 (Laplacian Solvers and Electrical Flows)}, we use expander decompositions to build near-linear time
        Laplacian system solvers.
\end{itemize}
\end{focusbox}

\section{Exercises}
\label{ch34:exercises}

\begin{exercise}[Warm-up cuts]
For the $4$-dimensional hypercube $Q_4$ ($16$ vertices, $4$-regular), compute $\abs{\partial S}/\abs S$ for $S=$ \{strings with first bit $0$\}, for $S=$ a single vertex, and for $S=$ a Hamming ball of radius $1$. Conjecture $h(Q_4)$ and argue the half-cube cut is optimal.
\end{exercise}

\begin{exercise}[The easy direction, your turn]
Prove $\lambda_2\le 2h(\Gph)$ for non-regular graphs using the \emph{normalized} Laplacian $\mathcal L=I-\Adjm_D^{-1/2}\Adjm\,\Adjm_D^{-1/2}$ and conductance $\phi$. Identify exactly where regularity was used in the proof of the easy direction of Theorem~\ref{ch34:thm-cheeger} and adapt it.
\end{exercise}

\begin{exercise}[Mixing Lemma is tight]
Show the Petersen graph nearly saturates Theorem~\ref{ch34:thm-eml}: find sets $S,T$ where $\abs{e(S,T)-\tfrac{d}{n}\abs S\abs T}$ is close to $\mu\sqrt{\abs S\abs T}$. (Hint: use eigenvectors of $\mu_2=1$ or $\mu_n=-2$ to choose $S,T$.)
\end{exercise}

\begin{exercise}[Bipartite poison]
Prove that a connected $d$-regular graph is bipartite iff $\mu_n=-d$. Conclude that bipartite graphs are never two-sided expanders, and explain why the lazy walk's $\tfrac12 I$ term is exactly the fix.
\end{exercise}

\begin{exercise}[Implement Cheeger rounding]
Implement Algorithm~\ref{ch34:alg-cheeger}. Test it on (a) a cycle $C_{1000}$, (b) a $\sqrt n\times\sqrt n$ grid, (c) a random $3$-regular graph. Report the returned ratio versus $\lambda_2/2$ and $\sqrt{2d\lambda_2}$, and confirm the quadratic gap is visible on the cycle.
\end{exercise}

\begin{exercise}[Alon--Boppana, baby case]
For $d=3$, verify $2\sqrt{d-1}=2\sqrt2\approx 2.83$ and check that the Petersen graph's $\mu=2<2.83$ does not contradict the bound. Explain why a single small graph can beat $2\sqrt{d-1}$ while no infinite family can. (Hint: Alon--Boppana is asymptotic; the $o(1)$ slack matters at small $n$.)
\end{exercise}

\begin{exercise}[Expander Chernoff in code]
Empirically test Theorem~\ref{ch34:thm-chernoff}. On a random $4$-regular graph mark a random $1/3$ of vertices ``bad''; run $10^4$ length-$k$ walks for $k=10,20,40$; plot the empirical tail of (fraction of steps in the bad set) against the independent-sampling Chernoff bound. Confirm the expander walk concentrates almost as tightly.
\end{exercise}

\begin{exercise}[Break a code]
Build a tiny expander (Tanner) code from the bipartite graph of Figure~\ref{ch34:fig-code} extended to $7$ variables and $3$ checks. Run the flip-decoder on a $1$-bit error and on a $2$-bit error. Construct an error pattern where greedy flipping gets stuck, and relate the failure to a set $S$ that \emph{fails} the lossless expansion condition $\abs{N(S)}\ge(1-\eps)d\abs S$.
\end{exercise}

\begin{exercise}[Zig-zag bookkeeping]
For the zig-zag product of $\Gph$ ($D$-regular on $N$ vertices) with $H$ ($d$-regular on $D$ vertices), prove the product has $ND$ vertices and is $d^2$-regular. Explain in one paragraph why squaring $\Gph$ before each zig-zag step is needed to keep the spectral gap from decaying to zero over the iteration.
\end{exercise}

\begin{exercise}[Collapse the walk]
Justify the distance-collapse used in Table~\ref{ch34:tab-mixing}: because the Petersen graph is distance-transitive, the lazy walk from a fixed start keeps equal mass on all vertices at the same distance, so the $10$-state chain reduces to a $3$-state chain on distances $\set{0,1,2}$. Write down that $3\times3$ transition matrix explicitly, find its eigenvalues, and check that its nontrivial ones are $\set{\tfrac12+\tfrac{1}{6},\,\tfrac12-\tfrac{2}{6}}=\set{2/3,\,1/6}$, matching $w_i=\tfrac12+\mu_i/(2d)$ for $\mu_i\in\set{1,-2}$. Why does the start-symmetry guarantee the collapse is exact rather than approximate?
\end{exercise}

\begin{exercise}[Resistance and expansion]
Using the electrical-network view (Chapter~35), argue that in a $d$-regular expander the effective resistance between \emph{any} two vertices is $\Th(1/d)$ up to lower-order terms, in sharp contrast to the cycle where it grows linearly. Connect this to commute times via the identity $\text{commute}(u,v)=2\abs E\cdot R_{\mathrm{eff}}(u,v)$.
\end{exercise}
